\section{The stack, the container, the PDK}

\begin{frame}{The full stack in one picture}
  \begin{center}
  \begin{tikzpicture}[node distance=0.4cm]
    \node[pipelinebox, minimum width=3.3cm, minimum height=1.0cm]
      (ll) {LibreLane\\\scriptsize Python orchestrator -- one YAML in};
    \node[pipelinebox2, below=1.2cm of ll, minimum width=1.8cm] (yo) {Yosys};
    \node[pipelinebox2, left=of yo,  minimum width=1.8cm] (or) {OpenROAD};
    \node[pipelinebox2, right=of yo, minimum width=1.8cm] (kl) {KLayout};
    \node[pipelinebox2, right=of kl, minimum width=1.8cm] (mg) {Magic};
    \node[pipelinebox2, right=of mg, minimum width=1.8cm] (ng) {Netgen};
    \draw[arrow] (ll) -- (or);
    \draw[arrow] (ll) -- (yo);
    \draw[arrow] (ll) -- (kl);
    \draw[arrow] (ll) -- (mg);
    \draw[arrow] (ll) -- (ng);
  \end{tikzpicture}
  \end{center}
  \vspace{0.4em}
  The CLI is literally \cmd{librelane config.yaml}. Everything below
  the top box is what does the \emph{actual} EDA work; LibreLane reads
  the config, sequences the tools, passes files between them, and
  writes one \cmd{metrics.csv}.
  \note{%
    This restates the thesis from section~1 before we zoom into the
    container and PDK. The message: there are \emph{six} independent
    open-source projects on this slide, not one. LibreLane, Yosys,
    OpenROAD, Magic, KLayout, Netgen are all maintained separately, each
    with its own release cadence and documentation.\par
    Name LibreLane as the evolution of OpenLane 2 (same lineage, Python
    rewrite). Audiences will almost never type \cmd{yosys ...} or
    \cmd{openroad ...} directly -- but when a step fails, the log you
    read is the underlying tool's log, not a LibreLane log.\par
    The upcoming two slides (container + PDK) are about giving these six
    tools a common place to live and a common set of foundry files to
    read.}
\end{frame}

\begin{frame}{Host, container, and your files}
  \begin{center}
  \begin{tikzpicture}[x=1cm, y=1cm]
    \draw[draw=muted, thick, rounded corners=3pt, fill=softbg]
      (-6.2, -2.2) rectangle (6.2, 2.0);
    \node[anchor=north, font=\small\bfseries, text=muted]
      at (0, 1.95) {Your Linux host};
    \draw[draw=accent, thick, rounded corners=3pt, fill=accent!6]
      (-6.0, -2.0) rectangle (2.2, 1.5);
    \node[anchor=north west, font=\scriptsize\bfseries, text=accent]
      at (-5.95, 1.45) {Docker container (\texttt{hpretl/iic-osic-tools:next})};
    \node[pipelinebox2, minimum width=7.6cm, minimum height=0.8cm,
          align=center, font=\scriptsize]
      at (-1.9, 0.35)
      {LibreLane + Yosys + OpenROAD + KLayout + Magic + Netgen};
    \node[pipelinebox3, minimum width=3.4cm, minimum height=1.0cm, font=\small]
      at (-3.7, -0.95) {PDK\\\scriptsize gf180mcuD};
    \node[pipelinebox, minimum width=3.4cm, minimum height=1.0cm, font=\small]
      at (-0.1, -0.95) {Your design\\\scriptsize (mounted)};
    \node[draw=accent2, thick, rounded corners=3pt, fill=accent2!10,
          minimum width=3.0cm, minimum height=1.4cm, align=center,
          font=\footnotesize] (hostdir)
      at (4.4, -0.95) {\texttt{\char126/eda/designs}\\\tiny (your files)};
    \draw[<->, thick, accent2] (1.8, -0.95) -- (2.9, -0.95);
    \node[font=\tiny, text=accent2, align=center] at (2.35, -0.5) {bind\\mount};
  \end{tikzpicture}
  \end{center}
  \vspace{0.2em}
  \small The container is the \term{environment}; the bind-mount is the
  \term{bridge} -- edits on the host appear instantly inside.
  \note{%
    Mental model is the whole point here. People comfortable with VMs get
    it in 30 seconds; people who have never used containers need to hear:
    the container is disposable, your host folder is not. Treat
    \cmd{\char126/eda/designs} as the canonical workspace. Everything
    written inside \cmd{/foss/designs} survives \cmd{docker rm}. Everything
    outside that path inside the container is gone.\par
    The bind-mount is the only reason we can run a heavy-weight image and
    still keep our editor, git, and artifacts on the host where we already
    have backups.}
\end{frame}

\begin{frame}[fragile]{Container bootstrap -- four commands}
  \begin{lstlisting}[style=shellstyle]
# 1. One-time image pull (~15 GB)
docker pull hpretl/iic-osic-tools:next

# 2. Start a long-running, headless container
docker run -d --name gf180 \
    -v ~/eda/designs:/foss/designs:rw \
    --user $(id -u):$(id -g) \
    hpretl/iic-osic-tools:next --skip sleep infinity

# 3. (Optional) interactive shell for debugging
docker exec -it gf180 bash

# 4. Non-interactive use (what the rest of the deck does)
docker exec gf180 bash -lc '<command>'
  \end{lstlisting}
  \note{%
    Four commands is all a first-time user needs. The flags matter:
    \cmd{-v} binds the workspace, \cmd{--user} keeps artifacts on the host
    owned by the right UID, \cmd{--skip sleep infinity} is the image's
    built-in ``do nothing, stay alive'' entrypoint so we can
    \cmd{docker exec} into it later.\par
    The \cmd{bash -lc} pattern is the non-interactive workhorse we use for
    every subsequent command; it yields a login shell so the container's
    \cmd{\$PATH}, \cmd{\$PDK\_ROOT}, etc. are set up. Use single quotes on
    the host side so the entire payload is handed to the container shell
    untouched.\par
    Cleanup (not on slide): \cmd{docker stop gf180 \&\& docker rm gf180}
    removes the container but leaves your host files untouched.}
\end{frame}

\begin{frame}[fragile]{GF180MCU: two PDK gotchas to know upfront}
  \begin{itemize}\small
    \item The image ships \cmd{gf180mcuD} (5\,V MCU flavour) via
          \term{ciel}. Good to go.
    \item The full-chip template needs two extra I/O cells
          (\cmd{gf180mcu\_ws\_io\_\_dvdd}, \cmd{\_\_dvss}) that live in the
          \term{wafer-space fork}:
          \begin{lstlisting}[style=shellstyle]
git clone --depth 1 --branch 1.8.0 \
  https://github.com/wafer-space/gf180mcu.git gf180mcu
          \end{lstlisting}
    \item The container defaults to IHP SG13G2. Every new shell must
          re-activate GF180:
          \begin{lstlisting}[style=shellstyle]
source sak-pdk-script.sh gf180mcuD gf180mcu_fd_sc_mcu7t5v0
          \end{lstlisting}
  \end{itemize}
  \note{%
    This is the single slide that saves every first-timer 2 hours of
    debugging. Without the wafer-space fork, you get an obscure ``unknown
    cell'' error during synthesis. Without the \cmd{sak-pdk-script}
    activation, LibreLane cheerfully looks for \cmd{sg13g2\_stdcell} files
    inside the GF180 directory and stops on a Tcl glob error.\par
    Both issues show up in Section 6 ``Pitfalls'' with their exact error
    messages -- mention that now so people recognise the symptoms later.
    Tag 1.8.0 of the wafer-space fork is the version we audited; bump only
    with intent.\par
    A \term{PDK} itself is just foundry files (stdcell libs, LEF, Liberty,
    SPICE models, DRC decks). We do not need to explain it in depth for a
    first-timer; the container already has the bits.}
\end{frame}
